\chapter{Системийн хэрэгжүүлэлт}

\section{Өгөгдлийн сангийн бүрдүүлэлт ба шошгололт}
Машин сургалтын загвар нь чанартай өгөгдлөөс шууд хамаардаг. Бид судалгааны хүрээнд бодит орчинд буюу дэлгүүрийн лангуун дээрх ундаа, ус, жүүс зэрэг төрөл бүрийн барааны зургийг гар утсаар цуглуулсан. Цуглуулсан $300$ гаруй зургуудыг өөр өөр өнцөг, алслалт, гэрэлтүүлэгтэйгээр авч баяжуулсан бөгөөд эдгээр зургууд дээрх объектуудыг \texttt{LabelImg} хэрэгслийг ашиглан гар аргаар шошголж (annotate), YOLO форматаар хадгалсан. Өгөгдлийн багцыг (Dataset) Сургалт (Train) $70\%$, Баталгаажуулалт (Validation) $20\%$, Туршилт (Test) $10\%$ гэсэн харьцаагаар хувааж загварын оролтод бэлтгэсэн. Түүнчлэн, сургалтын өгөгдлийг нэмэгдүүлэх зорилгоор мэдээллийг өсгөх (Data augmentation) аргууд болох эргүүлэх, өнгө өөрчлөх зэргийг ашигласан.

\section{Загвар сургалт (Model Training)}
Сургалтыг Python програмчлалын хэл болон PyTorch сангийн орчинд Ultralytics-ийн YOLOv8 API-г ашиглан хэрэгжүүлсэн. Загварын параметрүүд (hyperparameters)-ийг тохируулахдаа:
\begin{itemize}
    \item \textbf{Эпох (Epochs):} 50 удаагийн давтамжтай сургасан.
    \item \textbf{Багцын хэмжээ (Batch size):} 16
    \item \textbf{Зургийн хэмжээ (Image Size):} Оролтын зургийг 640x640 харьцаагаар тохируулж өгсөн.
\end{itemize}
Сургалтын үр дүнд хамгийн сайн нарийвчлалтай зогссон үеийн жинг (weights) \texttt{best.pt} файлаар хадгалан авч цаашдын Inference хөгжүүлэлтэд ашигласан.

\section{Серверийн хөгжүүлэлт (FastAPI Backend)}
Backend серверийг Python FastAPI ашиглан бүтээсэн. FastAPI нь асинхрон (\texttt{async/await}) дэмждэг тул сүлжээний хүсэлтийг хурдан боловсруулдаг давуу талтай. Мэдээллийг урт хугацаанд найдвартай хадгалах үүднээс NoSQL бүтэцтэй MongoDB өгөгдлийн санг Motor асинхрон драйвертай хоршуулан ашигласан \cite{fastapi}. Сервер нийтдээ гурван үндсэн үйлчилгээний үүргийг (Router Layer) гүйцэтгэнэ:
\begin{enumerate}
    \item \textbf{Products API}: Системд бүртгэлтэй бараануудын мэдээллийг CRUD хийх.
    \item \textbf{Detection API}: `multipart/form-data` хэлбэрээр зургийг хүлээн авч үйлчилгээний давхаргаар дамжуулан Machine Learning модель руу илгээж, танигдсан координатуудыг JSON хэлбэрээр буцаах.
    \item \textbf{Audit API}: Аудитын горим ажиллаж, танигдсан барааны жагсаалтад аудит хийж түүхийг хадгалах.
\end{enumerate}

\section{Мобайл хөгжүүлэлт (Flutter Frontend)}
Хэрэглэгчийн гар утасны аппликейшнийг Dart хэлэн дээр суурилсан Flutter фреймворк ашиглан хэрэгжүүлсэн. Аппликейшн дотор \texttt{camera} болон \texttt{image\_picker} сангуудыг ашиглаж шууд зураг авч, \texttt{http} сангаар дамжуулан REST API руу POST хүсэлт илгээнэ. Мэдээллийн төлвийг удирдахдаа \texttt{Provider} ашигласан ба илрүүлэлтийн болон аудитын үр дүнг (Pass, Fail, Warning) онцлох өнгөнүүдийн хамтаар (UI components) хэрэглэгчдэд ойлгомжтой байдлаар дүрслэн харуулсан. Зургийг байршуулах хугацаанд Loading дэлгэц болон алдаа гарсан үед алдааг мэдээлэх механизмыг бүрэн суулгасан болно \cite{flutter}.
